\scanpage{10}
\uline{Ex}\quad\tikz[remember picture,baseline]\node[inner sep=0pt] (afftitle) {Groupe affine \unsure{associé} à un groupe $\unsure{\mathbb{Z}|Z}$};

\diagram{\begin{tikzcd}[ampersand replacement=\&,column sep=small,row sep=1.2em]
\& \& \operatorname{Aff}(\mathbb{Z})\subset\unsure{\mathfrak{S}|\mathcal{S}}_{\mathbb{Z}_{d}}
   \arrow[d,phantom,"{\rotatebox{18}{$=$}}"{description,xshift=-4.5em}] \& \& \\
1 \arrow[r] \& \mathbb{Z} \arrow[r,"i"] \&
\operatorname{Aut}_{\substack{\text{ens à}\\\text{gr d'op}}}
 (\mathbb{Z},\mathbb{Z}_{d}) \arrow[r,"p"] \arrow[d,hook] \&
\operatorname{Aut}(\mathbb{Z}) \arrow[r] \arrow[l,bend left=29,"q"] \& 1\\
\& \& \operatorname{Aut}(\mathbb{Z})\times\mathfrak{S}_{\mathbb{Z}_{d}} \& \&
\end{tikzcd}}

\noindent\begin{minipage}[t]{.50\linewidth}
\[
\begin{aligned}
&\gamma\in\operatorname{Aut}(\mathbb{Z})\\
&\tikz[remember picture,baseline=(s010qline.base)]{\node[inner sep=0pt] (s010qline) {$q(\gamma)=(\gamma,\gamma)$};}\\
&i(g)=(\mathrm{id}_{\mathbb{Z}},\tau_{g^{\unsure{\prime|.|}}}^{\unsure{s|\deleted{s}|\delta}})
  \tikz[remember picture,baseline]\coordinate (idefend);\\
&\boxed{\operatorname{int}(q(\gamma))\mid\mathbb{Z}=\gamma}
\end{aligned}
\]
\end{minipage}\hfill
\begin{minipage}[t]{.45\linewidth}
\diagram{\begin{tikzcd}[ampersand replacement=\&,row sep=.1em]
\underbrace{(\gamma,\rho)\in\operatorname{Aut}(\mathbb{Z},\mathbb{Z}_{d})}\\
\Updownarrow\\
\rho(gx)=\gamma(g)\rho(x)\\
\forall g,x\in\mathbb{Z}
\end{tikzcd}}
\end{minipage}
% Detached upper-gap point and lower i(g) point; only the lower point has uncertain status as writing.
\figure{\begin{tikzpicture}[remember picture,overlay]
\fill ([xshift=5.1em,yshift=.3ex]s010qline.east) circle (.45pt);
\node[inner sep=0pt] at ([xshift=-.15em,yshift=-.65em]idefend) {$\unsure{.}$};
\end{tikzpicture}}

\[
\operatorname{Aff}_{\deleted{\unsure{\substack{\text{ens à}\\\text{gr d'op}}}}}
 (\mathbb{Z})\deleted{\unsure{(\mathbb{Z}_{d})|(G_d)}}\simeq
 \underset{\text{\tikz[baseline=-.4ex]{\draw (0,.04)--(.13,0);}}}{\mathbb{Z}}\mathbin{.}\operatorname{Aut}(\mathbb{Z})
 \qquad\text{\unsure{(produit|produit|(produits|produits} $1/2$ direct)}
\]
% The short oblique stroke below the first product factor is detached; no mathematical role is assigned to it.

Si $\Gamma\subset\operatorname{Aut}(\mathbb{Z})$ ($\Gamma$ ss-groupe) on a un ss-groupe naturel de $\deleted{\unsure{\operatorname{Aut}_{\substack{\text{ens}\\\text{gr d'op}}}(G,G_d)}}\operatorname{Aff}(\mathbb{Z})\deleted{\unsure{=|-}}$, savoir $\operatorname{Aff}^{\Gamma}(\mathbb{Z})\simeq\mathbb{Z}\mathbin{.}\Gamma$.

\noindent\tikz[remember picture,baseline]\coordinate (affblocktop);
Hom d'ens à op. avec \unsure{groupe|groupes} \unsure{variable|variables}
\diagram{\begin{tikzcd}[ampersand replacement=\&,column sep=large]
(E,\unsure{\mathbb{G}|G})\arrow[r,"{(\rho,\gamma)}"']\&(E',\unsure{G'|C'})
\end{tikzcd}}

\deleted{\uline{NB}}

\uline{NB} Si $E'$ est un torseur sous \tikz[baseline=(s010target.base)]{\node[inner sep=0pt] (s010target) {$G'\vphantom{\rule{0pt}{3.8ex}}$};\node[overlay,anchor=base west,inner sep=0pt] (s010nonempty) at ([xshift=.6em,yshift=2.6ex]s010target.base east) {\inserted{et $E\ne\varnothing$}};\draw[overlay] ([xshift=.1em,yshift=-.2ex]s010target.base east) .. controls +(.2em,2ex) and +(1em,0) .. ([xshift=.2em]s010nonempty.south east);}, alors $\rho$ détermine $\gamma$ car \deleted{\unsure{pour}} si $x\in E$, $\forall g\in G$, $\gamma(g)$ est déterminé en termes de $\rho$ par la condition $\rho(g.x)=\gamma(g).\rho(x)$
\tikz[remember picture,baseline]\coordinate (affblockbottom);
\figure{\begin{tikzpicture}[remember picture,overlay]
\draw ([xshift=-.3em]affblocktop.west) .. controls +(-.7em,0) and +(-.25em,.65em) .. ([xshift=-1em,yshift=-.55em]affblocktop.west)
 -- ([xshift=-1em,yshift=.7em]affblocktop.west |- affblockbottom)
 .. controls +(0,-.55em) and +(-.3em,0) .. ([xshift=-.35em]affblocktop.west |- affblockbottom);
\draw[->] ([xshift=-1em,yshift=-.65em]affblocktop.west)
 .. controls +(-3em,2em) and +(-4em,-1em) .. ([xshift=-2em,yshift=-.45em]afftitle.west)
 -- ([yshift=-.45em]afftitle.west) -- ([yshift=.65em]afftitle.west);
\end{tikzpicture}}

Cas $\mathbb{Z}$ \tikz[baseline=(s010vectoriel.base)]{\node[inner sep=0pt] (s010vectoriel) {vectoriel\vphantom{\rule{0pt}{3.8ex}}};\node[overlay,anchor=base west,inner sep=0pt] (s010underlying) at ([xshift=.3em,yshift=2.6ex]s010vectoriel.base west) {\inserted{grpe \unsure{ss-jacent à|s-jacent à|sous-jacent à}}};\draw[overlay] ([xshift=-.2em,yshift=-.2ex]s010vectoriel.base west) .. controls +(-.1em,2ex) and +(1em,0) .. ([xshift=.2em]s010underlying.south east);} $V$ sur corps $R$ (réels p ex)
\[
\begin{aligned}
&\Gamma=\operatorname{Aut}_{R}(V)\subset\operatorname{Aut}_{\mathrm{gr}}(V)\\[.8em]
&\operatorname{Aff}_{R}(V)\simeq V\mathbin{.}\operatorname{Aut}_{R}(V)
\end{aligned}
\]
\[
(\underset{\substack{\rotatebox{90}{$\ni$}\\V}}{z},
 \underset{\substack{\rotatebox{90}{$\ni$}\\\underset{R}{\operatorname{Aut}(V)}}}{\gamma})
 \mathbin{.}(z',\gamma')=(z+\gamma(z'),\gamma\gamma')
\]

Si $V$ muni d'une forme quadratique non \tikz[baseline=(s010deg.base)]{\node[inner sep=0pt] (s010deg) {dég\unsure{,|.}};\node[anchor=south,inner sep=0pt] at ([yshift=1ex]s010deg.north east) {\inserted{$Q$}};} \unsure{on|$\to$} a le groupe des déplacements (transf. affines qui \tikz[remember picture,baseline]\node[inner sep=0pt] (distanceverb) {\unsure{invariant|invarient}}; la distance)
% A separate short stroke stands beneath the opening of the distance verb; its purpose is uncertain.
\figure{\begin{tikzpicture}[remember picture,overlay]
\draw ([xshift=.55em,yshift=-.22em]distanceverb.south west) -- ++(.65em,.02em);
\end{tikzpicture}}
\[
\operatorname{Depl}_{\deleted{\unsure{\text{aff}|\text{Aff}|\text{orth}}}}(V,Q)\simeq V\mathbin{.}\operatorname{Orth}(V,Q)
\]
% archive sheet number: 9
